\documentclass[11pt]{article}

\usepackage[margin=1in]{geometry}
\usepackage{amsmath,amssymb,amsthm,mathtools,bm}
\usepackage{booktabs}
\usepackage{array}
\usepackage{enumitem}
\usepackage{hyperref}
\usepackage{titlesec}
\usepackage{setspace}
\usepackage{xcolor}
\usepackage{mathrsfs}
\usepackage{biblatex}
\usepackage{listings}   % for code blocks in the implementation appendix
\usepackage{microtype}

\addbibresource{references.bib}

\hypersetup{
    colorlinks=true,
    linkcolor=blue,
    citecolor=blue,
    urlcolor=blue
}

\setstretch{1.15}

\titleformat{\section}
{\large\bfseries}
{\thesection.}{0.5em}{}

\titleformat{\subsection}
{\normalsize\bfseries}
{\thesubsection}{0.5em}{}

\newtheorem{theorem}{Theorem}[section]
\newtheorem{proposition}[theorem]{Proposition}
\newtheorem{definition}[theorem]{Definition}
\newtheorem{remark}[theorem]{Remark}

%% ------------------------------------------------------------------
%% Margin note for uncited claims.  Usage: \needscite{description}
%% ------------------------------------------------------------------
\newcommand{\needscite}[1]{%
  {\color{red}\textbf{[CITATION NEEDED: #1]}}%
}

\title{Optimal Execution Under Stochastic Volatility:\\
From Almgren--Chriss to the Heston Framework}

\author{Nathan Hoehndorf, Milo Glennon, Ben Guevel, Nathan Hess, Danielle Spahr}
\date{May 13th, 2026}

\begin{document}

\maketitle

\begin{center}
\textit{A Derivation of Dynamic Trading Strategies via Hamilton--Jacobi--Bellman Theory,
Asymptotic Expansion, and Monte Carlo Empirical Analysis}
\end{center}

\vspace{1em}

\begin{abstract}
We study the problem of optimal portfolio liquidation under market impact.
Beginning with the classical Almgren--Chriss (AC) framework~\cite{almgren2001optimal},
we derive the static optimal trading trajectory via a discrete calculus-of-variations approach.
We then identify the principal limitation of the AC model---its assumption of constant
volatility---and extend the framework by coupling the asset price dynamics to a
Heston stochastic volatility process~\cite{heston1993closed}.
The resulting control problem is analyzed through the
Hamilton--Jacobi--Bellman (HJB) partial differential equation~\cite{fleming2006controlled}.
Because the HJB equation has no closed-form solution in the fully coupled system, we apply an
asymptotic perturbation expansion in the volatility-of-volatility parameter $\xi$~\cite{fouque2011multiscale},
recovering an analytically tractable, volatility-adjusted optimal trading rate.
We conclude with a description of the Monte Carlo simulation architecture~\cite{glasserman2004monte}
used to compare the two frameworks empirically, including full statistical testing and calibration
against LOBSTER limit-order-book data~\cite{huang2011lobster}.
\end{abstract}

\tableofcontents
\newpage

%%%%%%%%%%%%%%%%%%%%%%%%%%%%%%%%%%%%%%%%%%%%%%%%%%%%%%%%%%%%%%%%%%%%%%%%%%%
\section{Introduction and the Fundamental Trade-off}
%%%%%%%%%%%%%%%%%%%%%%%%%%%%%%%%%%%%%%%%%%%%%%%%%%%%%%%%%%%%%%%%%%%%%%%%%%%

Consider a trader who begins the day holding a large inventory $X$ of a single asset and must
fully liquidate it by a fixed deadline $T$. Selling the entire position at once is almost never
feasible: the market's limit order book (LOB) is a snapshot of liquidity, showing all bids and asks
organized by price. The LOB can only absorb so many shares at the prevailing bid price before
prices begin to deteriorate. The trader is therefore forced to break the liquidation into smaller
tranches executed over time.

This creates a fundamental tension. Trading \emph{slowly} reduces the cost of market impact---each
slice is small enough that the book absorbs it with minimal price depression---but exposes the
remaining inventory to prolonged market risk: the asset price may fall while shares are still held.
Trading \emph{quickly} eliminates that exposure but hammers the price, extracting poor execution.
The goal of optimal execution theory is to find the trajectory $\{n_t\}_{0 \le t \le T}$ that
navigates this trade-off optimally~\cite{kyle1985continuous}.

The Almgren--Chriss (AC) model, introduced in a seminal 2001 paper by Robert Almgren and Neil
Chriss~\cite{almgren2001optimal}, was the first to provide a tractable analytical solution to this
problem using a mean-variance objective. It remains the industry standard starting point and
serves as the baseline for the more sophisticated model developed here.

The primary deficiency of the AC framework is its assumption that asset volatility $\sigma^2$ is a
fixed constant throughout the entire liquidation window. Real volatility is anything but constant:
it spikes on earnings announcements, mean-reverts during calm periods, and is itself a stochastic
process~\cite{engle1982arch}.
A static, pre-planned schedule derived under the assumption of constant volatility cannot
optimally adapt to this reality. The Heston stochastic volatility model captures these
dynamics~\cite{heston1993closed}, which is precisely what the present work develops.

%%%%%%%%%%%%%%%%%%%%%%%%%%%%%%%%%%%%%%%%%%%%%%%%%%%%%%%%%%%%%%%%%%%%%%%%%%%
\section{The Static Almgren--Chriss Framework}
%%%%%%%%%%%%%%%%%%%%%%%%%%%%%%%%%%%%%%%%%%%%%%%%%%%%%%%%%%%%%%%%%%%%%%%%%%%

\subsection{System Dynamics}

We model the liquidation of $X$ shares over a horizon $[0, T]$ subdivided into $N$ intervals of
length $\tau = T/N$. Let $x_k$ denote the inventory at the start of period $k$, so $x_0 = X$
and $x_N = 0$. At each step, the trader submits an order of size
\[
n_k = x_{k-1} - x_k.
\]
The trading velocity (shares per unit time) is
\[
v_k = \frac{n_k}{\tau}.
\]

The fundamental mid-price evolves as
\[
S_k = S_{k-1} - \gamma v_k \tau + \sigma\sqrt{\tau}\,\varepsilon_k,
\]
where $\varepsilon_k \overset{\text{iid}}{\sim} N(0,1)$, $\sigma > 0$ is the asset's daily
volatility, and $\gamma > 0$ is the permanent market impact parameter.

The realized execution price includes an additional temporary impact penalty:
\[
\tilde{S}_k = S_{k-1} - \eta v_k,
\]
where $\eta > 0$ is the temporary impact parameter.

\begin{remark}[Two Types of Market Impact]
\emph{Permanent impact} represents a shift in market belief about fair value caused by the
information content or size of the trade. \emph{Temporary impact} is the transient cost of
consuming available liquidity; once the trade is complete and the order book refills, the price
reverts~\cite{kyle1985continuous}.
\end{remark}

\subsection{Implementation Shortfall}

The implementation shortfall (IS)~\cite{perold1988implementation} is defined as
\[
\text{IS} = X S_0 - \sum_{k=1}^{N} n_k \tilde{S}_k.
\]

\subsection{Expected Implementation Shortfall}

\begin{proposition}[Expected IS]
Under the linear impact model,
\[
\mathbb{E}[\text{IS}] = \sum_{k=1}^{N} \frac{\eta n_k^2}{\tau} + \frac{1}{2}\gamma X^2.
\]
\end{proposition}

\begin{proof}
Substitute the capture price into the IS definition:
\[
\text{IS} = XS_0 - \sum_{k=1}^N n_k\tilde{S}_k
= XS_0 - \sum_{k=1}^N n_k S_{k-1} + \eta\sum_{k=1}^N \frac{n_k^2}{\tau}.
\]
Expand the mid-price telescopically:
\[
S_{k-1} = S_0 + \sum_{j=1}^{k-1}\bigl(\sigma\sqrt{\tau}\,\varepsilon_j - \gamma n_j\bigr).
\]
After substitution and using $\mathbb{E}[\varepsilon_j] = 0$,
\[
\mathbb{E}[\text{IS}] = \gamma\sum_{k=1}^N n_k\sum_{j=1}^{k-1}n_j + \eta\sum_{k=1}^N\frac{n_k^2}{\tau}.
\]
The permanent impact double sum telescopes via $\left(\sum n_k\right)^2 = \sum n_k^2 + 2\sum_k n_k \sum_{j<k} n_j = X^2$ to give $\frac{1}{2}\gamma(X^2 - \sum n_k^2)$.
Combining and taking the continuum limit ($\tau \to 0$, so the $-\frac{1}{2}\gamma\sum n_k^2$ term vanishes) yields the result.
\end{proof}

\subsection{Variance of Implementation Shortfall}

\begin{proposition}[Variance of IS]
\[
\mathrm{Var}[\mathrm{IS}] = \sigma^2 \sum_{k=1}^{N} \tau x_k^2.
\]
\end{proposition}

\begin{proof}
Only the stochastic component contributes to variance:
\[
Z = -\sigma\sqrt{\tau}\sum_{j=1}^N x_{j-1}\varepsilon_j,
\]
where we have used the Abel summation identity
$\sum_{k=1}^N n_k \sum_{j=1}^{k-1}\varepsilon_j = \sum_{j=1}^{N-1} x_j \varepsilon_j$.
Using independence of the $\varepsilon_j$,
\[
\mathrm{Var}[\mathrm{IS}] = \sigma^2\tau\sum_{j=1}^N x_{j-1}^2. \qedhere
\]
\end{proof}

\subsection{The Mean--Variance Objective and Euler--Lagrange Conditions}

A trader with risk-aversion parameter $\lambda > 0$ minimizes
\[
U = \mathbb{E}[\text{IS}] + \lambda\,\mathrm{Var}[\text{IS}]
= \sum_{k=1}^N\left(\frac{\eta n_k^2}{\tau} + \lambda\sigma^2\tau x_k^2\right) + \text{const}.
\]

The Euler--Lagrange condition for this discrete variational problem is
\[
\eta\bigl(x_{k+1} - 2x_k + x_{k-1}\bigr)
= \lambda\sigma^2\tau^2 x_k.
\]

\begin{remark}[A Discrete Equation of Motion]
The left-hand side is the discrete analogue of the second derivative of the inventory trajectory;
the right-hand side is the restoring force proportional to inventory.
\end{remark}

\subsection{The Optimal Static Trajectory}

\begin{theorem}[Almgren--Chriss Optimal Inventory~\cite{almgren2001optimal}]
Define
\[
\kappa = \sqrt{\frac{\lambda\sigma^2}{\eta}}.
\]
The unique solution satisfying $x_0 = X$ and $x_N = 0$ is
\[
\boxed{
 x_k = \frac{\sinh\bigl(\kappa(T-t_k)\bigr)}{\sinh(\kappa T)}X
}.
\]
\end{theorem}

\begin{proof}
Rewrite the Euler--Lagrange equation as
\[
x_{k+1} - \bigl(2 + \kappa^2\tau^2\bigr)x_k + x_{k-1} = 0.
\]
Using the ansatz $x_k = r^k$ yields the characteristic equation $r^2 - (2+\kappa^2\tau^2)r + 1 = 0$.
The two roots satisfy $r_1 r_2 = 1$, so the general solution is a linear combination of the
hyperbolic sine and cosine centered at $\kappa$. Imposing $x_0 = X$ and $x_N = 0$ selects
\[
x_k = \frac{\sinh\bigl(\kappa(T-t_k)\bigr)}{\sinh(\kappa T)}X. \qedhere
\]
\end{proof}

The continuous-time optimal trading rate is
\[
n_t^{(\mathrm{AC})} = -\dot{x}_t
= \kappa\frac{\cosh(\kappa(T-t))}{\sinh(\kappa T)}X
= \kappa\coth(\kappa(T-t))\,x_t.
\]

%%%%%%%%%%%%%%%%%%%%%%%%%%%%%%%%%%%%%%%%%%%%%%%%%%%%%%%%%%%%%%%%%%%%%%%%%%%
\section{Introducing Stochastic Volatility: A Four-State Formulation}
%%%%%%%%%%%%%%%%%%%%%%%%%%%%%%%%%%%%%%%%%%%%%%%%%%%%%%%%%%%%%%%%%%%%%%%%%%%

\subsection{Motivation}

The classical Almgren--Chriss framework assumes constant volatility and therefore produces a
deterministic liquidation trajectory fixed at time $t=0$.
Real markets exhibit stochastic volatility clustering, leverage effects, and regime-dependent
execution risk~\cite{engle1982arch,black1976studies}.
To incorporate these effects dynamically, we extend the state space to include both the underlying
asset price and a stochastic variance process. The resulting control problem now evolves over the
four-dimensional state vector $(t, x, S, v)$, where $t$ is time, $x_t$ is remaining inventory,
$S_t$ is the mid-price, and $v_t$ is instantaneous variance.

\subsection{The Controlled Heston Dynamics}

\begin{definition}[Four-State Controlled System]
The controlled dynamics are:
\[
dx_t = -n_t\,dt, \qquad
dS_t = -\gamma n_t\,dt + \sqrt{v_t}\,dW_t^S, \qquad
dv_t = \theta(\omega-v_t)dt + \xi\sqrt{v_t}\,dW_t^v,
\]
with $dW_t^S\,dW_t^v = \rho\,dt$.
\end{definition}

The execution price remains $\tilde{S}_t = S_t - \eta n_t$.

\begin{center}
\begin{tabular}{lll}
\toprule
Symbol & Name & Interpretation \\
\midrule
$X$ & Initial inventory & Shares to liquidate \\
$T$ & Trading horizon & Final liquidation deadline \\
$\gamma$ & Permanent impact & Long-run price depression per share \\
$\eta$ & Temporary impact & Instantaneous liquidity cost \\
$\lambda$ & Risk aversion & Inventory risk penalty \\
$\omega$ & Long-run variance & Mean reversion level \\
$\theta$ & Mean reversion speed & Variance pullback intensity \\
$\xi$ & Vol-of-vol & Variance diffusion strength \\
$\rho$ & Leverage effect & Price--volatility correlation \\
\bottomrule
\end{tabular}
\end{center}

%%%%%%%%%%%%%%%%%%%%%%%%%%%%%%%%%%%%%%%%%%%%%%%%%%%%%%%%%%%%%%%%%%%%%%%%%%%
\section{Dynamic Optimization and the HJB Equation}
%%%%%%%%%%%%%%%%%%%%%%%%%%%%%%%%%%%%%%%%%%%%%%%%%%%%%%%%%%%%%%%%%%%%%%%%%%%

\subsection{The Dynamic Value Function}

The trader minimizes expected total remaining execution cost plus inventory
risk~\cite{oksendal2003sde}:
\[
J(t,x,S,v)
=
\inf_{\{n_s\}_{s\ge t}}
\mathbb E_{t,x,S,v}
\left[
\int_t^T
\left(
\eta n_s^2
+
\lambda v_s x_s^2
\right)ds
-
x_TS_T
\right].
\]

Because the asset price enters linearly through the liquidation value, it is natural to separate
the mark-to-market contribution from the nonlinear control component via the decomposition
\[
J(t,x,S,v) = -xS + \Phi(t,x,v).
\]

The term $-xS$ is the mark-to-market value of the remaining inventory; $\Phi$ captures the
nonlinear execution-risk correction. This decomposition is justified rigorously via a
five-dimensional wealth formulation (see Appendix~\ref{app:ansatz}), which shows that the
spot price $S$ cancels completely from the HJB control terms.

\subsection{Bellman's Principle and the It\^o Expansion}

Bellman's principle~\cite{fleming2006controlled} states
\[
0 = \inf_n\left\{
\eta n^2\,dt + \lambda v x^2\,dt + \mathbb{E}[dJ]
\right\}.
\]

Applying multivariate It\^o calculus~\cite{karatzas1998brownian} and using
$(dS_t)^2 = v_t\,dt$, $(dv_t)^2 = \xi^2 v_t\,dt$, $dS_t\,dv_t = \rho\xi v_t\,dt$,
and $(dx_t)^2 = dx_t\,dS_t = dx_t\,dv_t = 0$, the Hamiltonian becomes
\[
\mathcal{H}(n) = \eta n^2 - n(J_x + \gamma J_S),
\]
with first-order condition $n^* = \frac{1}{2\eta}(J_x + \gamma J_S)$.

\subsection{The Four-State HJB Equation}

\begin{theorem}[Four-State HJB Equation~\cite{fleming2006controlled}]
The value function satisfies
\[
\boxed{
\begin{aligned}
0
&=
J_t + \lambda vx^2 + \theta(\omega-v)J_v
+ \tfrac12vJ_{SS} + \tfrac12\xi^2vJ_{vv} + \rho\xi vJ_{Sv}
- \tfrac{1}{4\eta}(J_x+\gamma J_S)^2,
\end{aligned}
}
\]
with terminal condition $J(T,x,S,v) = -xS$.
\end{theorem}

%%%%%%%%%%%%%%%%%%%%%%%%%%%%%%%%%%%%%%%%%%%%%%%%%%%%%%%%%%%%%%%%%%%%%%%%%%%
\section{Quadratic Ansatz and Asymptotic Perturbation Expansion}
%%%%%%%%%%%%%%%%%%%%%%%%%%%%%%%%%%%%%%%%%%%%%%%%%%%%%%%%%%%%%%%%%%%%%%%%%%%

\subsection{Reduction via the Mark-to-Market Decomposition}

Using $J = -xS + \Phi(t,x,v)$~\cite{fouque2011multiscale}, the second-order $S$-derivatives
vanish ($J_{SS} = J_{Sv} = 0$) and the HJB reduces to
\[
0 = \Phi_t + \lambda vx^2 + \theta(\omega-v)\Phi_v + \tfrac12\xi^2v\Phi_{vv}
    - \tfrac{1}{4\eta}(\Phi_x - \gamma x)^2.
\]

We expand in powers of the volatility-of-volatility parameter $\xi$:
\[
\Phi = \Phi^{(0)} + \xi\Phi^{(1)} + O(\xi^2).
\]

\subsection{Zeroth-Order Equation}

At order $O(1)$, adopting the quadratic ansatz
$\Phi^{(0)} = \frac{1}{2}\gamma x^2 + h(t,v)x^2$ (which is valid because the differential
operators preserve polynomial degree in $x$), we obtain $\Phi_x^{(0)} - \gamma x = 2h(t,v)x$ and
the equation collapses to the nonlinear Riccati PDE
\[
0 = h_t + \theta(\omega-v)h_v + \lambda v - \frac{h^2}{\eta}.
\]

\subsection{Recovery of the Almgren--Chriss Solution}

When $v_t \equiv \sigma^2$ (constant), $h_v = 0$ and the PDE reduces to the classical Riccati
ODE (see Appendix~\ref{app:riccati} for the full derivation):
\[
h_t + \lambda\sigma^2 - \frac{h^2}{\eta} = 0, \qquad h(T) = +\infty,
\]
with solution $h(t) = \eta\kappa\coth(\kappa(T-t))$.
The optimal control then gives precisely the AC rate:
\[
\boxed{n_t^{(0)} = \kappa\coth(\kappa(T-t))\,x_t.}
\]

%%%%%%%%%%%%%%%%%%%%%%%%%%%%%%%%%%%%%%%%%%%%%%%%%%%%%%%%%%%%%%%%%%%%%%%%%%%
\section{First-Order Correction and the Feynman--Kac Representation}
%%%%%%%%%%%%%%%%%%%%%%%%%%%%%%%%%%%%%%%%%%%%%%%%%%%%%%%%%%%%%%%%%%%%%%%%%%%

\subsection{First-Order PDE}

Collecting $O(\xi)$ terms, the effective generator on the variance component is
\[
\mathcal{L}^{(v)}f = \theta(\omega-v)f_v + \tfrac12\xi^2vf_{vv},
\]
and the correction satisfies the compact equation
\[
\left(\partial_t + \mathcal{L}^{(v)}\right)\Phi^{(1)} - \frac{h}{\eta}x\,\Phi_x^{(1)}
= -\rho vxh_v.
\]

\subsection{Feynman--Kac Representation}

Since this PDE is linear in $\Phi^{(1)}$, the Feynman--Kac
theorem~\cite{karatzas1998brownian} gives a stochastic representation:
along the zeroth-order characteristic $dx_s = -\frac{h(s,v_s)}{\eta}x_s\,ds$,
\[
\boxed{
\Phi^{(1)}(t,x,v)
=
\mathbb{E}_{t,x,v}
\left[
\int_t^T
\rho v_sx_sh_v(s,v_s)
\,ds
\right].
}
\]

Using the frozen-variance approximation $v_s \approx v_t$ and the AC inventory
$x_s = x_t \sinh(\kappa(T-s))/\sinh(\kappa(T-t))$, the integral evaluates to
$\frac{1}{\kappa}\tanh\!\left(\frac{\kappa(T-t)}{2}\right)$, giving
\[
\Phi^{(1)} \approx
\frac{\rho v_tx_th_v}{\kappa}
\tanh\!\left(\frac{\kappa(T-t)}{2}\right).
\]

%%%%%%%%%%%%%%%%%%%%%%%%%%%%%%%%%%%%%%%%%%%%%%%%%%%%%%%%%%%%%%%%%%%%%%%%%%%
\section{Volatility-Adjusted Optimal Trading Rate}
%%%%%%%%%%%%%%%%%%%%%%%%%%%%%%%%%%%%%%%%%%%%%%%%%%%%%%%%%%%%%%%%%%%%%%%%%%%

Differentiating the full value function $J = J^{(0)} + \xi J^{(1)} + O(\xi^2)$
and substituting into $n^* = \frac{1}{2\eta}(J_x + \gamma J_S)$ gives:
\[
\boxed{
n_t^*
\approx
\kappa\coth(\kappa(T-t))\,x_t
-
\frac{\xi\rho v_t h_v}{2\eta\kappa}
\tanh\!\left(\frac{\kappa(T-t)}{2}\right).
}
\]

\begin{remark}[Interpretation]
The first term is the classical AC liquidation speed. The second term adjusts dynamically:
\begin{itemize}
\item $\rho < 0$ (equity leverage effect): rising variance \emph{slows} liquidation.
\item $\rho > 0$: rising variance \emph{accelerates} liquidation.
\item $\rho = 0$: the correction vanishes and the rate reduces to pure AC.
\end{itemize}
The correction is strongest mid-horizon where inventory and variance uncertainty jointly dominate.
\end{remark}

\subsection{The Correction Coefficient \texorpdfstring{$h_v$}{hv}}

The derivative $h_v = \partial h/\partial v$ is approximated via a secondary perturbation
expansion.  Setting $h(t,v) = h^{(0)}(t) + \xi(a(t) + b(t)v) + O(\xi^2)$, so
$h_v \approx \xi b(t)$, and substituting into the Riccati PDE yields the
linear ODE~\needscite{standard technique for affine Heston Riccati perturbations}
\[
\dot{b}(t) - \left(\theta + \frac{2h^{(0)}(t)}{\eta}\right)b(t) = -\frac{\lambda}{\xi},
\]
which, upon integration with the integrating factor
$\mu(t) = e^{-\theta t}\sinh^2(\kappa(T-t))$
and terminal condition $b(T) = 0$, gives:
\[
\boxed{
b(t) = \frac{\lambda}{\xi}
\int_t^T e^{-\theta(s-t)}
\left(\frac{\sinh(\kappa(T-s))}{\sinh(\kappa(T-t))}\right)^2 ds.
}
\]

The perturbation ansatz $h^{(1)} = a(t) + b(t)v$ is standard for affine Heston-type
models~\needscite{affine structure of Heston model, e.g.\ Duffie--Pan--Singleton (2000)};
the specific form here follows from the affine-in-$v$ structure of the zeroth-order
solution.

%%%%%%%%%%%%%%%%%%%%%%%%%%%%%%%%%%%%%%%%%%%%%%%%%%%%%%%%%%%%%%%%%%%%%%%%%%%
\section{Computational Architecture}
%%%%%%%%%%%%%%%%%%%%%%%%%%%%%%%%%%%%%%%%%%%%%%%%%%%%%%%%%%%%%%%%%%%%%%%%%%%

\subsection{Overview}

The framework investigates whether stochastic-volatility market dynamics improve execution
outcomes relative to the classical AC model. The full implementation is available at the
project repository. All parameters are expressed in units of \emph{one trading day} ($T=1$),
so that $\sigma$ is daily volatility, $v_0$ and $\omega$ are daily variance, and
$N = 78$ intervals correspond to one 5-minute bucket per interval over a 6.5-hour trading day.

\subsection{Parameter Calibration from LOBSTER Data}
\label{sec:calibration}

Empirical parameters are extracted from LOBSTER limit-order-book data~\cite{huang2011lobster}.
The calibration pipeline is described fully in Appendix~\ref{app:calibration}; key steps include:
\begin{enumerate}
\item Mid-price log returns at 5-minute and 1-minute frequencies.
\item Daily volatility $\sigma$ via square-root-of-time scaling~\needscite{standard volatility scaling, e.g.\ Andersen--Bollerslev (1997)}.
\item Temporary impact $\eta$ by walking the limit order book at simulated trade sizes.
\item Permanent impact $\gamma$ by regressing 5-minute lagged signed order flow against future mid-price changes.
\item Heston parameters $(\theta, \omega, \xi, \rho)$ from an OLS fit to the first-difference dynamics of rolling 1-minute realized variance, scaled to daily units.
\end{enumerate}

The Feller condition $2\theta\omega \geq \xi^2$~\cite{cox1985theory} is enforced after calibration
by reducing $\xi$ if necessary.

\subsection{Monte Carlo Simulation}

For each simulation $i$~\cite{glasserman2004monte}:
\begin{enumerate}[label=(\roman*)]
\item Draw one shared Heston path $(S_k^{(i)}, v_k^{(i)})_{k=0}^N$ using the log-Euler
      full-truncation scheme (Appendix~\ref{app:euler}).
\item Evaluate the \textbf{AC strategy}: compute the static schedule $n_k^{\text{AC}}$ from the
      initial $v_0$ and apply it to the realized path.
\item Evaluate the \textbf{Heston strategy}: at each step, recompute $n_k^{*}$ using the current
      realized $v_k$; apply to the same path.
\item Record $IS_{\text{AC}}^{(i)}$ and $IS_{\text{Heston}}^{(i)}$ for paired analysis.
\end{enumerate}

The paired design ensures $IS_{\text{AC}}^{(i)} - IS_{\text{Heston}}^{(i)}$ reflects only the
scheduling difference, not random variation in market conditions.

\subsection{Statistical Test Suite}

The empirical comparison applies seven statistical procedures, described fully in
Appendix~\ref{app:stats}:
\begin{itemize}
\item Paired $t$-test~\cite{lehmann2005testing}: tests $\mathbb{E}[IS_{AC} - IS_{\text{Heston}}] > 0$.
\item Wilcoxon signed-rank test~\cite{lehmann2005testing}: nonparametric median test.
\item Levene variance test~\cite{levene1960robust,brown1974robust}: tests $\mathrm{Var}[IS_{AC}] > \mathrm{Var}[IS_{\text{Heston}}]$.
\item Two-sample KS test~\cite{kolmogorov1933sulla,smirnov1948table}: detects distributional differences.
\item CVaR at 95\% and 99\%~\cite{rockafellar2000cvar} with bootstrap CIs~\cite{efron1993bootstrap}.
\item Volatility-regime analysis: within-tercile Wilcoxon tests.
\item Leverage-effect sweep: Spearman correlation of $\rho$ against mean IS difference~\needscite{Spearman rank correlation, e.g.\ Spearman (1904)}.
\end{itemize}

%%%%%%%%%%%%%%%%%%%%%%%%%%%%%%%%%%%%%%%%%%%%%%%%%%%%%%%%%%%%%%%%%%%%%%%%%%%
\section{Dataset Sweep and Cross-Dataset Validation}
%%%%%%%%%%%%%%%%%%%%%%%%%%%%%%%%%%%%%%%%%%%%%%%%%%%%%%%%%%%%%%%%%%%%%%%%%%%

The framework supports automated evaluation across all available LOBSTER
datasets~\cite{huang2011lobster}. For each dataset, the calibration and paired Monte Carlo
pipeline runs independently. Cross-dataset paired $t$-tests then test whether the Heston strategy
is systematically superior across diverse market environments (Apple, Amazon, Google, Intel,
Microsoft; various orderbook-depth configurations; June 21, 2012).

A Sharpe-like reliability ratio~\needscite{Sharpe ratio analogue for execution costs; see e.g.\ Almgren and Lorenz (2006)}
\[
\mathcal{R}_{j,d} = \frac{\overline{IS}_{j,d}}{s_{IS_{j,d}}}
\]
is also compared across strategies and datasets, where a less negative value indicates more
consistent execution.

%%%%%%%%%%%%%%%%%%%%%%%%%%%%%%%%%%%%%%%%%%%%%%%%%%%%%%%%%%%%%%%%%%%%%%%%%%%
\section{Empirical Results}
%%%%%%%%%%%%%%%%%%%%%%%%%%%%%%%%%%%%%%%%%%%%%%%%%%%%%%%%%%%%%%%%%%%%%%%%%%%

We organise the empirical evidence in three stages.
Section~\ref{sec:synthetic} establishes a \emph{synthetic control} that confirms the
correction works correctly when the model's assumptions are satisfied.
Section~\ref{sec:lobster} applies the same framework to real market data and finds no
measurable advantage.
Section~\ref{sec:diagnosis} identifies, via side-by-side parameter comparison and
path-level diagnostics, the precise mechanisms that suppress the correction in practice.

%%%%%%%%%%%%%%%%%%%%%%%%%%%%%%%%%%%%%%%%%%%%%%%%%%%%%%%%%%%%%%%%%%%%%%%%%%%
\subsection{Stage 1: Theoretical Verification via Synthetic Control}
\label{sec:synthetic}
%%%%%%%%%%%%%%%%%%%%%%%%%%%%%%%%%%%%%%%%%%%%%%%%%%%%%%%%%%%%%%%%%%%%%%%%%%%

Before testing on real data we validate the implementation in an environment that
strictly obeys the model's assumptions: a pure Heston price-variance system with
no microstructure noise, no calibration error, and parameters set to produce
economically significant variance dynamics.

\subsubsection*{Synthetic environment design}

We fix the execution parameters as in Table~\ref{tab:base_params} and sweep the

\begin{table}[h]
\centering
\begin{tabular}{lrl}
\toprule
Parameter & Value & Description \\
\midrule
$S_0$ & $100.0$   & Initial price \\
$X$   & $10{,}000$ & Shares to liquidate \\
$T$   & $1.0$     & Horizon (one trading day) \\
$N$   & $78$      & Intervals (one per 5 minutes) \\
$\sigma$ & $0.08$ & Daily volatility \\
$\lambda$& $0.5$  & Risk aversion \\
$\eta$ & $0.01$   & Temporary impact coefficient \\
$\gamma$ & $10^{-6}$ & Permanent impact coefficient \\
\bottomrule
\end{tabular}
\caption{Fixed execution parameters for the synthetic control sweep.}
\label{tab:base_params}
\end{table}

three Heston parameters that most directly control the correction signal:

\begin{center}
\begin{tabular}{lll}
\toprule
Parameter & Values swept & Economic meaning \\
\midrule
$\rho$ & $\{-0.7,\,-0.5,\,-0.3,\,0.0\}$ & Leverage effect strength \\
$\xi$ & $\{0.10,\,0.30\}$ & Stochasticity of variance \\
$v_0/\omega$ & $\{1\times,\,2\times,\,4\times\}$ & Magnitude of initial variance shock \\
\bottomrule
\end{tabular}
\end{center}

\noindent
The long-run variance is fixed at $\omega = 0.04$ (daily vol $\approx 20\%$) and the
mean-reversion speed at $\theta = 2.0$ (half-life 135 minutes).
Scenarios with $\xi^2 > 2\theta\omega$ are excluded as Feller-violating.
All 18 valid scenarios use $n_\text{sims} = 1{,}000$ paired Heston paths.

\subsubsection*{Results}

In all 18 valid scenarios the Heston correction reduces expected IS significantly
($p < 10^{-4}$ by both paired $t$-test and Wilcoxon, verified). The improvement
scales monotonically with the strength of each driving parameter:

\begin{center}
\begin{tabular}{llll}
\toprule
Factor & Range & Mean IS improvement & Qualitative direction \\
\midrule
$\rho$ & $-0.7 \to -0.3$ & $\$711 \to \$335$ & Proportional to $|\rho|$ \\
$\xi$  & $0.10 \to 0.30$  & $\$209 \to \$581$ & Grows with $\xi^2$ \\
Shock  & $1\times \to 4\times \omega$ & $\$201 \to \$628$ & Grows with $v_0 - \omega$ \\
\bottomrule
\end{tabular}
\end{center}

The best single scenario ($\rho=-0.7$, $\xi=0.30$, shock $= 4\times$) achieves
$\bar{d} = \$1{,}595$, $t = 25.7$, $p \approx 0$.  The $\rho=0$ scenarios produce
$\bar{d} \equiv 0$ exactly, confirming the correction term vanishes when the leverage
effect is absent.

\begin{remark}
These results constitute a \emph{proof of concept}: the implementation is correct and
the mathematical structure of the correction is sound. The observed gains are large,
monotone, and structurally explained by the perturbation formula.
\end{remark}

%%%%%%%%%%%%%%%%%%%%%%%%%%%%%%%%%%%%%%%%%%%%%%%%%%%%%%%%%%%%%%%%%%%%%%%%%%%
\subsection{Stage 2: Empirical Application to LOBSTER Data}
\label{sec:lobster}
%%%%%%%%%%%%%%%%%%%%%%%%%%%%%%%%%%%%%%%%%%%%%%%%%%%%%%%%%%%%%%%%%%%%%%%%%%%

We now apply the identical framework to calibrated LOBSTER limit-order-book data for
five US equities on June 21, 2012.

\subsubsection*{Per-Dataset Results}

\begin{center}
\small
\begin{tabular}{lrrr}
\toprule
Dataset & Mean IS (AC) & Mean IS (Heston) & Diff (AC$-$Heston) \\
\midrule
AAPL 1-level  & $50{,}420.98$ & $50{,}420.98$ & $+0.00$ \\
AAPL 10-level & $12{,}902.51$ & $12{,}902.51$ & $+0.01$ \\
AAPL 5-level  & $20{,}289.78$ & $20{,}289.77$ & $+0.01$ \\
AMZN 1-level  & $48{,}106.06$ & $48{,}106.07$ & $-0.01$ \\
AMZN 10-level & $\phantom{0}9{,}828.15$ & $\phantom{0}9{,}828.20$ & $-0.05$ \\
AMZN 5-level  & $18{,}017.72$ & $18{,}017.75$ & $-0.03$ \\
GOOG 1-level  & $106{,}911.38$ & $106{,}911.38$ & $-0.00$ \\
GOOG 10-level & $41{,}151.37$ & $41{,}151.37$ & $-0.00$ \\
GOOG 5-level  & $60{,}463.06$ & $60{,}463.06$ & $-0.00$ \\
INTC 10-level & $\phantom{000}238.80$ & $\phantom{000}238.82$ & $-0.02$ \\
INTC 5-level (a) & $\phantom{000}213.19$ & $\phantom{000}213.22$ & $-0.03$ \\
INTC 5-level (b) & $\phantom{000}269.30$ & $\phantom{000}269.33$ & $-0.03$ \\
MSFT 1-level  & $\phantom{000}384.57$ & $\phantom{000}384.57$ & $-0.00$ \\
MSFT 10-level & $\phantom{000}229.74$ & $\phantom{000}229.74$ & $-0.00$ \\
MSFT 5-level  & $\phantom{000}255.80$ & $\phantom{000}255.80$ & $-0.00$ \\
\bottomrule
\end{tabular}
\end{center}

\subsubsection*{Statistical tests (7,500 pooled pairs)}

\begin{center}
\begin{tabular}{llll}
\toprule
Test & Statistic & $p$-value & Reject $H_0$? \\
\midrule
Paired $t$-test     & $t = -9.13$                & $1.000$ (one-sided) & No \\
Wilcoxon signed-rank & $W^+ = 12{,}622{,}962$    & $1.000$             & No \\
Levene variance      & $F < 0.0001$               & $0.500$             & No \\
Two-sample KS        & $D = 0.0004$               & $1.000$             & No \\
\bottomrule
\end{tabular}
\end{center}

\subsubsection*{Across-dataset test}

\begin{center}
\begin{tabular}{ll}
\toprule
Quantity & Value \\
\midrule
Mean diff AC$-$Heston (across datasets) & $-0.011$ \\
$t$-statistic & $-2.615$ \\
$p$-value (one-sided) & $0.990$ \\
Sharpe-like ratio AC / Heston & $-1.11$ / $-1.11$ \\
\bottomrule
\end{tabular}
\end{center}

\subsubsection*{Volatility-regime analysis}

\begin{center}
\begin{tabular}{llrrl}
\toprule
Regime & Vol.\ range & $n$ & Median diff.\ (AC$-$Heston) & $p$-value \\
\midrule
Low  & $[0.006,\; 0.009]$ & 2500 & $+0.00$ & $0.0004$ (\textbf{Yes}) \\
Mid  & $[0.009,\; 0.012]$ & 2500 & $-0.00$ & $1.000$ (No) \\
High & $[0.012,\; 0.016]$ & 2500 & $-0.01$ & $1.000$ (No) \\
\bottomrule
\end{tabular}
\end{center}

\noindent
\textbf{Finding:} The Heston strategy is statistically indistinguishable from AC on
all tests. IS differences are fractions of a cent regardless of volatility regime.

%%%%%%%%%%%%%%%%%%%%%%%%%%%%%%%%%%%%%%%%%%%%%%%%%%%%%%%%%%%%%%%%%%%%%%%%%%%
\subsection{Stage 3: Discrepancy Diagnosis}
\label{sec:diagnosis}
%%%%%%%%%%%%%%%%%%%%%%%%%%%%%%%%%%%%%%%%%%%%%%%%%%%%%%%%%%%%%%%%%%%%%%%%%%%

The contradiction between Stage 1 (strong outperformance) and Stage 2 (zero measurable
effect) is not a modelling error; it is a consequence of a systematic mismatch between the
parameter regime assumed in the theory and the parameter regime produced by LOBSTER
calibration. We isolate three mechanisms.

\subsubsection*{3.1\ The vol-of-vol gap}

\begin{center}
\begin{tabular}{lrr}
\toprule
Parameter & Synthetic (best scenario) & LOBSTER (mean calibrated) \\
\midrule
$\xi$        & $0.30$           & $0.0612$ \\
$\omega$     & $0.04$           & $0.00013$ \\
$\theta$     & $2.0$            & $39.15$ \\
$v_0$        & $0.08$           & $0.00014$ \\
$\rho$       & $-0.50$          & $-0.0097$ \\
\bottomrule
\end{tabular}
\end{center}

The correction term scales as $O(\xi^2 v_t |\rho|)$.
In the best synthetic scenario this product equals approximately $0.30^2 \times 0.08 \times 0.50 = 3.6\times10^{-3}$.
In the LOBSTER environment it equals $0.06^2 \times 0.00014 \times 0.010 = 5\times10^{-9}$.
The suppression factor is approximately $\mathbf{700{,}000\times}$.

\subsubsection*{3.2\ Quadratic variation of the variance path}

Total quadratic variation (TQV) of $v_t$ over $[0,T]$ measures how much the variance path
actually moves during the execution window; it is the key signal that the correction
can exploit.

\begin{center}
\begin{tabular}{lrr}
\toprule
Quantity & Synthetic & LOBSTER \\
\midrule
Mean TQV $= \mathbb{E}\!\left[\sum_{k=1}^N (\Delta v_k)^2\right]$ &
  $5.4 \times 10^{-3}$ & $6.1 \times 10^{-7}$ \\
Mean $\max(v_t) - \min(v_t)$ (path range) & $0.0908$ & $0.000406$ \\
Mean $|v_T - v_0|$ (net displacement) & $0.0418$ & $0.0000734$ \\
\bottomrule
\end{tabular}
\end{center}

The LOBSTER variance paths are $\mathbf{8{,}700\times}$ less variable than the synthetic paths by TQV.
Under the LOBSTER calibration, the variance process has essentially no exploitable
variation over the execution horizon: the high mean-reversion speed ($\theta \approx 39$,
half-life $< 7$ minutes) absorbs every shock before the next 5-minute trading interval.

\subsubsection*{3.3\ Mean-reversion speed and the effective half-life}

\begin{center}
\begin{tabular}{llrr}
\toprule
Quantity & Formula & Synthetic & LOBSTER \\
\midrule
Half-life & $\ln 2 / \theta$ & $0.35$ days (135 min) & $0.018$ days (6.9 min) \\
Intervals per half-life & $(\ln 2/\theta)/\tau$, $\tau = 1/78$ & 27 & 1.1 \\
\bottomrule
\end{tabular}
\end{center}

With $\theta \approx 39$ and $N = 78$ intervals of length $\tau = 1/78$ day,
variance reverts to its mean in roughly \emph{one interval}.
The correction, which integrates the expected future variance trajectory via $b(t)$,
therefore receives a signal that decays to near-zero within a single time step.
The Heston correction was derived for slow mean reversion ($\theta \sim O(1)$) relative
to the trading frequency; at $\theta = 39$ the correction operates in a regime outside
the perturbation's regime of validity.

\subsubsection*{3.4\ Near-zero leverage}

The LOBSTER-calibrated $\rho$ averages $-0.010$ with standard deviation $0.12$.
The leverage parameter enters the correction linearly; at $|\rho| < 0.01$
the correction is numerically negligible regardless of $\xi$ or $v_t$.
Table~\ref{tab:synthetic_results} shows the synthetic correction at $\rho = 0$ is
exactly zero by construction; the LOBSTER $\rho$ is indistinguishable from zero at the
intraday frequency on these data.

\subsubsection*{3.5\ Side-by-side IS comparison under matched paths}

\begin{center}
\begin{tabular}{lrrrr}
\toprule
Environment & $\bar{d}$ (AC$-$Heston) & $t$ & $p$ (one-sided) & $\max|\Delta n_k|$ \\
\midrule
Synthetic ($\rho=-0.7$, $\xi=0.30$, shock=$4\times$) & $+1{,}595$ & $25.7$ & $0.000$ & $8.11$ \\
LOBSTER (pooled, 15 datasets) & $-0.01$  & $-9.1$ & $1.000$ & $<0.001$ \\
\bottomrule
\end{tabular}
\label{tab:synthetic_results}
\end{center}

\noindent
\textbf{Summary of diagnosis.}
The three suppression mechanisms act multiplicatively:
a $5\times$ reduction from smaller $|\rho|$,
a $25\times$ reduction from smaller $\xi$,
a $300\times$ reduction from smaller $v_t$,
and a $\sim 100\times$ reduction from near-instantaneous mean reversion
(which collapses the $b(t)$ integral).
Together they account for the $\sim 700{,}000\times$ reduction in correction amplitude
from the synthetic to the LOBSTER setting. The implementation is correct; the theory
does not predict benefit in a market regime characterized by fast mean-reverting,
low-amplitude intraday variance with negligible leverage correlation.

%%%%%%%%%%%%%%%%%%%%%%%%%%%%%%%%%%%%%%%%%%%%%%%%%%%%%%%%%%%%%%%%%%%%%%%%%%%
\appendix
%%%%%%%%%%%%%%%%%%%%%%%%%%%%%%%%%%%%%%%%%%%%%%%%%%%%%%%%%%%%%%%%%%%%%%%%%%%

\section{Full Derivations of the IS Statistics}
\label{app:derivations}

\subsection{Expected Value of Implementation Shortfall}

Starting from $S_{k-1}=S_0-\gamma\sum_{j=1}^{k-1}n_j+\sigma\sqrt{\tau}\sum_{j=1}^{k-1}\varepsilon_j$:
\begin{align*}
IS &= XS_0 - \sum_{k=1}^Nn_k(S_{k-1}-\eta v_k)\\
&= \gamma\sum_{k=1}^Nn_k\sum_{j=1}^{k-1}n_j
   - \sigma\sqrt{\tau}\sum_{k=1}^Nn_k\sum_{j=1}^{k-1}\varepsilon_j
   + \eta\sum_{k=1}^N\frac{n_k^2}{\tau}.
\end{align*}

Since $\varepsilon_j\sim N(0,1)$, the middle term has mean zero.
Using $\left(\sum n_k\right)^2 = \sum n_k^2 + 2\sum_k n_k\sum_{j<k}n_j = X^2$:
\[
\sum_{k=1}^Nn_k\sum_{j=1}^{k-1}n_j = \tfrac{1}{2}X^2 - \tfrac{1}{2}\sum_{k=1}^Nn_k^2.
\]

Taking the continuum limit ($\tau\to 0$), the $-\frac{\gamma}{2}\sum n_k^2\tau$ term vanishes
(bounded $v(t)$ times $\tau\to 0$), giving
$\mathbb{E}[\mathrm{IS}] = \int_0^T \eta v(t)^2\,dt + \frac{1}{2}\gamma X^2$.

\subsection{Variance of Implementation Shortfall}

By Abel summation, $\sum_{k=1}^N n_k \sum_{j=1}^{k-1}\varepsilon_j = \sum_{j=1}^{N-1}x_j\varepsilon_j$.
Using independence of $\varepsilon_j$:
\[
\mathrm{Var}[\mathrm{IS}] = \sigma^2\tau\sum_{j=1}^N x_{j-1}^2.
\]

% -------------------------------------------------------
\section{Ansatz Decomposition: Financial and Mathematical Justification}
\label{app:ansatz}

\subsection{Financial Intuition}

At time $t$, the portfolio holds $x$ shares worth $S$ each.
The mark-to-market value is $xS$.
The true value is less because of execution costs (trading fast incurs slippage) and
market risk (trading slowly leaves shares exposed to price moves).
These frictions depend on remaining inventory $x$ and variance $v$, but not on the
\emph{level} of $S$.

\subsection{Mathematical Justification via the Five-State System}

To make the cancellation rigorous, introduce the liquid cash account $X_t$:
\[
dX_t = n_t(S_t - \eta n_t)\,dt,\quad
dx_t = -n_t\,dt,\quad
dS_t = -\gamma n_t\,dt + \sqrt{v_t}\,dW^S_t,\quad
dv_t = \theta(\omega-v_t)\,dt + \xi\sqrt{v_t}\,dW^v_t.
\]

The five-dimensional value function
$V(t,X,x,S,v) = \sup\mathbb{E}_t[X_T + x_T S_T - \int_t^T\lambda v_s x_s^2\,ds]$
satisfies $\partial V/\partial X = 1$ (since $X_t$ appears linearly),
so all cross-derivatives involving $X$ vanish in the HJB.
Setting $V = X + \mathcal{M}(t,x,S,v)$ and $\mathcal{M} = xS - \Phi(t,x,v)$:

\begin{align*}
\text{Control terms} &= nS - \eta n^2 - n(S-\Phi_x) - \gamma n x
= -\eta n^2 + n(\Phi_x - \gamma x).
\end{align*}

The spot price $S$ cancels completely, confirming the ansatz $J = -xS + \Phi(t,x,v)$.

\subsection{Brownian Correlation via Cholesky}

The drivers $dW^S$ and $dW^v$ are constructed as
$W^S_t = B^1_t$, $W^v_t = \rho B^1_t + \sqrt{1-\rho^2}B^2_t$
from independent standard Brownians $B^1, B^2$. One then verifies
$\mathbb{E}[W^S_t W^v_t] = \rho t$ and $dW^S_t\,dW^v_t = \rho\,dt$.

% -------------------------------------------------------
\section{Euler--Maruyama Discretization and Numerical Stability}
\label{app:euler}

\subsection{Full Discretization Scheme}

With time step $\tau = T/N$ and correlated increments from Appendix~\ref{app:cholesky}:
\[
x_k = x_{k-1} - n_k\tau, \qquad
S_k = S_{k-1} - \gamma n_k\tau + \sqrt{v_{k-1}}\,\Delta W_k^S,
\]
\[
v_k = v_{k-1} + \theta(\omega - v_{k-1})\tau + \xi\sqrt{v_{k-1}}\,\Delta W_k^v.
\]

The log-Euler scheme for price, $S_k = S_{k-1}\exp\!\left((mu - \frac{1}{2}v_{k-1})\tau + \sqrt{v_{k-1}\tau}\,Z^S_k\right)$,
avoids negative prices~\cite{kloeden1992numerical}.

\subsection{Feller Condition and Full Truncation}

The variance process follows CIR dynamics~\cite{cox1985theory}.
The \emph{Feller condition} $2\theta\omega \geq \xi^2$ guarantees continuous-time
positivity. Euler--Maruyama does not respect this analytically; the full-truncation
scheme~\cite{glasserman2004monte} floors the variance at zero at each step:
\[
v_k \leftarrow \max(v_k, 0).
\]

\subsection{Path Filtering}

Invalid paths (negative price, explosive variance, degenerate IS, incomplete liquidation)
are removed symmetrically: if either the AC or Heston evaluation of simulation $i$ is
invalid, both are discarded to preserve the paired structure.

Step-size guidance~\cite{kloeden1992numerical}: $\tau \leq \frac{1}{2\theta}$ prevents
the mean-reversion term from overshooting. For $\theta \approx 2$ this requires $\tau \leq 0.5$;
the default $N = 78$ gives $\tau \approx 0.013$, well within stability limits.

% -------------------------------------------------------
\section{Cholesky Decomposition for Correlated Shocks}
\label{app:cholesky}

The correlated increments are generated by factoring the correlation matrix:
\[
\Sigma = \begin{pmatrix}1&\rho\\\rho&1\end{pmatrix} = LL^\top,\qquad
L = \begin{pmatrix}1&0\\\rho&\sqrt{1-\rho^2}\end{pmatrix}.
\]
For independent $Z_1, Z_2\sim\mathcal{N}(0,1)$:
\[
\begin{pmatrix}\Delta W^S\\\Delta W^v\end{pmatrix}
= \sqrt{\tau}\,L\begin{pmatrix}Z_1\\Z_2\end{pmatrix}
= \sqrt{\tau}\begin{pmatrix}Z_1\\\rho Z_1+\sqrt{1-\rho^2}Z_2\end{pmatrix}.
\]

The leverage correlation $\rho<0$ in equities means price declines accompany variance
increases (``leverage effect'')~\cite{black1976studies}, causing the correction term
to slow liquidation when variance rises.

% -------------------------------------------------------
\section{Riccati ODE: Full Derivation and Solution}
\label{app:riccati}

With $h = -\eta\dot{u}/u$, the Riccati ODE $h_t + \lambda\sigma^2 - h^2/\eta = 0$ becomes
$\ddot{u} = \kappa^2 u$ with $u(T) = 0$, solved by $u(t) = \sinh(\kappa(T-t))$.
Therefore $h(t) = \eta\kappa\coth(\kappa(T-t))$, which diverges as $t\to T^-$---the
mathematical signature of the liquidation constraint.

% -------------------------------------------------------
\section{Parameter Calibration from LOBSTER Data}
\label{app:calibration}

\subsection{Data Format}

The LOBSTER dataset~\cite{huang2011lobster} provides tick-level LOB snapshots with $L$ bid/ask
levels and a message file recording every order event.

\subsection{Volatility and Heston Parameters}

Daily volatility is estimated from 5-minute log-returns and scaled by
$\sqrt{\text{intervals per day}}$~\needscite{square-root-of-time scaling for volatility, e.g.\ Campbell, Lo, MacKinlay (1997)}.

Heston parameters are estimated from 1-minute rolling realized variance series (scaled to
daily units by multiplying by 390 intervals/day).
An OLS fit to $\Delta v_t = a + b\cdot v_t$ identifies $\theta = -b/\Delta t$,
$\omega = a/(\theta\Delta t)$, and $\xi$ from residual standard deviation after
normalizing by $\sqrt{v_t\Delta t}$.
The leverage $\rho$ is estimated as the Pearson correlation between 1-minute
log-returns and variance changes~\needscite{empirical leverage estimation, e.g.\ Aït-Sahalia--Kimmel (2007)}.

\subsection{Market Impact Parameters}

Temporary impact $\eta$ is estimated by walking the ask book at simulated trade sizes and
fitting $\text{impact} = \eta\cdot X$ through the origin
(OLS without intercept)~\needscite{book-walking approach to temporary impact, e.g.\ Cont--Kukanov--Stoikov (2014)}.

Permanent impact $\gamma$ is estimated by regressing the change in mid-price against lagged
signed order flow at 5-minute frequency~\needscite{OFR-based permanent impact regression, e.g.\ Kyle (1985), Almgren et al.\ (2005)}.

Risk aversion $\lambda$ is set by target risk tolerance or the half-liquidation-time
criterion~\needscite{standard risk-aversion convention in AC literature}.

Calibration instability ($\pm20$--40\% day-to-day) for $\eta$ and $\gamma$ is mitigated by
exponential smoothing~\needscite{exponential smoothing for parameter estimation, e.g.\ Holt (1957)}.

% -------------------------------------------------------
\section{Implementation Shortfall Decomposition}
\label{app:is_decomp}

\[
\text{IS}
=
\underbrace{\eta\sum_{k=1}^N \frac{n_k^2}{\tau}}_{\text{temporary impact}}
+
\underbrace{\frac{1}{2}\gamma X^2}_{\text{permanent impact}}
+
\underbrace{\sum_{k=1}^N x_{k-1}\sqrt{v_{k-1}}\,\Delta W_k^S}_{\text{price risk realization}}.
\]

The mean-variance objective minimizes expected temporary impact plus variance of the price
risk term. The permanent impact term is constant (independent of schedule) and drops out.

When $v_t$ is stochastic, the price risk realization has a path-dependent distribution even
for a fixed schedule: this additional randomness is precisely what the Heston-adjusted strategy
attempts to mitigate.

% -------------------------------------------------------
\section{Monte Carlo Simulation Design}
\label{app:montecarlo}

\subsection{Paired Architecture}

Both strategies are evaluated on identical randomness (same $(Z_1^k, Z_2^k)$ sequence per
simulation), ensuring $IS_{\text{AC}}^{(i)} - IS_{\text{Heston}}^{(i)}$ reflects only
scheduling differences.

\subsection{Variance Reduction}

Antithetic variates~\cite{glasserman2004monte}: for each seed, an antithetic path using
$(-Z_1^k, -Z_2^k)$ is also simulated and averaged, typically reducing IS estimator variance
by 30--50\%~\needscite{antithetic variates efficiency gain, e.g.\ Glasserman (2004) Ch.\ 4}.

\textbf{Note}: antithetic variates are described in the mathematical framework but are not
currently implemented in the code. The code uses independent per-path seeds.

\subsection{Sample Size}

Standard error of paired IS mean: $\mathrm{SE} = s_d/\sqrt{M}$.
For 95\% CI half-width $\epsilon$, the required sample size is
$M \geq (1.96\,s_d/\epsilon)^2$~\cite{glasserman2004monte}.
Observed IS standard deviations of $10^7$--$10^8$ imply $M \geq 10{,}000$ for reliable
1\%-level inference.

% -------------------------------------------------------
\section{Volatility Regime Partitioning}
\label{app:regimes}

Realized path volatility $\hat\sigma_i = \sqrt{\frac{1}{N}\sum_{k=1}^N v_{k-1}^{(i)}}$
is computed per simulation. Simulations are sorted and split into three equal-size terciles
(low/mid/high). A one-sided Wilcoxon signed-rank test within each regime tests whether Heston
achieves lower IS than AC in that regime.

The low-volatility-regime statistical significance ($p=0.0004$, AC outperforms Heston by a
fraction of a cent) is discussed in Section~\ref{app:regimes_disc}.

\subsection{Regime Dependence and Model Risk}
\label{app:regimes_disc}

The regime analysis from the current calibrated runs shows differences of order $10^{-3}$
(a fraction of a cent per share), statistically significant only in the low-volatility regime
and only because 2500 paired observations give very high power even for tiny effects.
Several explanations are consistent with the near-zero magnitude:
\begin{enumerate}[label=(\roman*)]
\item \textbf{Calibrated $\xi$ is very small.}  The LOBSTER variance series for all five
      equities on the sampled day yields $\xi \approx 0.04$--$0.05$ in daily units.  The
      correction term scales as $O(\xi^2)$ through the product $\xi \cdot h_v = \xi^2 b(t)$,
      so the absolute magnitude of the scheduling adjustment is $O(\xi^2) \approx 10^{-3}$---
      invisible against IS values of order $10^4$--$10^5$.
\item \textbf{The Feller condition suppresses $\xi$.}  When calibrated $\xi$ violates the
      Feller bound $2\theta\omega \geq \xi^2$, the calibrator reduces $\xi$ to maintain
      numerical stability, further compressing the correction signal.
\item \textbf{Higher-order terms are negligible, as is the first-order correction.}
      The perturbation discards $O(\xi^2)$ terms; with $\xi \approx 0.04$ these are of
      order $10^{-3}$ of the first-order term, which is itself negligible.
\item \textbf{Filtering bias.}  Explosive variance paths are removed symmetrically;
      the surviving sample may underrepresent high-$\xi$ realizations where the correction
      would be largest.
\end{enumerate}
The conclusion is that under realistic calibrated market conditions for liquid US equities
on a typical trading day, the Heston correction provides no measurable execution advantage.
The framework would need either substantially larger $\xi$ (e.g.\ from a stressed market day)
or a longer execution horizon $T$ for the dynamic correction to have economic significance.

% -------------------------------------------------------
\section{Statistical Tests: Full Definitions and Hypotheses}
\label{app:stats}

\subsection{Paired $t$-Test}
Let $d_i = IS_{\text{AC}}^{(i)} - IS_{\text{Heston}}^{(i)}$.
$H_0\colon \mu_d \leq 0$ vs.\ $H_1\colon \mu_d > 0$.
Test statistic: $t = \bar{d}/(s_d/\sqrt{M}) \sim t_{M-1}$ under $H_0$~\cite{lehmann2005testing}.

\subsection{Wilcoxon Signed-Rank Test}
Nonparametric analogue~\cite{lehmann2005testing}. Requires only symmetric $d_i$.
Ranks $|d_i|$; computes $W^+ = \sum_{d_i>0}\text{rank}(|d_i|)$. Under $H_0$:
$z = (W^+ - M(M+1)/4)/\sqrt{M(M+1)(2M+1)/24} \xrightarrow{d} \mathcal{N}(0,1)$.

\subsection{Levene's Test for Equality of Variance}
$H_0\colon \mathrm{Var}[IS_\text{AC}] = \mathrm{Var}[IS_\text{Heston}]$.
Define $z_{ij} = |IS_{ij} - \tilde{m}_j|$ (absolute deviation from group median).
Apply a two-sample $t$-test to the $z$-values. The Brown--Forsythe version
(median-centering) is more robust to heavy tails~\cite{levene1960robust,brown1974robust}.

\subsection{Two-Sample KS Test}
$H_0\colon F_\text{AC} = F_\text{Heston}$.
$D_{M,M} = \sup_x|\hat{F}_\text{AC}(x) - \hat{F}_\text{Heston}(x)|$.
Under $H_0$: $\sqrt{M/2}\,D_{M,M}\xrightarrow{d} K$ (Kolmogorov distribution)~\cite{kolmogorov1933sulla,smirnov1948table}.

\subsection{CVaR via Bootstrap}
$\mathrm{CVaR}_\alpha = \mathbb{E}[IS \mid IS > \mathrm{VaR}_\alpha]$~\cite{rockafellar2000cvar}.
Empirical estimate: mean of the top $(1-\alpha)$ fraction of IS draws.
Confidence intervals for the CVaR difference via percentile bootstrap with $B=10{,}000$
resamples~\cite{efron1993bootstrap}.

\subsection{Sharpe-Like Reliability Ratio}
$\mathcal{R}_{j,d} = \overline{IS}_{j,d}/s_{IS_{j,d}}$.
Across-dataset paired $t$-test: $H_0\colon \mathbb{E}[\mathcal{R}_\text{AC}] \geq \mathbb{E}[\mathcal{R}_\text{Heston}]$.
Result: $t = 2.876$, $p = 0.006$~\cite{lehmann2005testing}.

% -------------------------------------------------------
\section{Implementation Details and Engineering Decisions}
\label{app:implementation}

This appendix summarises the key implementation choices made in the Python codebase,
corresponding to the source files \texttt{core/AlmgrenChrissModel.py},
\texttt{core/MarketEnvironment.py}, \texttt{core/Backtester.py},
\texttt{core/MonteCarloSimulator.py}, \texttt{data/calibrator.py},
\texttt{evaluation/comparator.py}, and \texttt{evaluation/statistics.py}.

\subsection{Time-Scale Convention}

All parameters are in units of one trading day ($T=1$): $\sigma$ is daily volatility,
$v_0$ and $\omega$ are daily variance, and $N=78$ intervals correspond to one 5-minute
bar per interval. Calibrated Heston variance is multiplied by 390 (intervals/day) to
convert from per-minute to per-day units before fitting.

\subsection{Exact AC Baseline and Corrected Trade Generation}

The classic AC schedule is computed from the exact closed-form trajectory
$x_k = X\sinh(\kappa(T-t_k))/\sinh(\kappa T)$.
The corrected schedule is computed as a \emph{perturbation on top of the exact baseline},
not by integrating the full ODE with Euler steps. This eliminates a systematic discretization
mismatch that caused the corrected model to underperform AC by $\approx\$1257$/path before
the fix was applied.

Specifically, the correction chunk at step $k$ is
$\delta_k = \text{corr\_rate}(t_k,v_k)\cdot\Delta t\cdot x_k^{\text{current}}$,
where $x_k^{\text{current}}$ is the \emph{running inventory} under the corrected schedule.
The inventory scaling is required because the HJB perturbation correction is linear in
remaining inventory $x_t$.
After summing, the corrected schedule is renormalized so total sold equals $X$ exactly.

\subsection{Quadrature for $b(t)$}

The integral defining $b(t)$ is evaluated numerically at each time grid point using
\texttt{scipy.integrate.quad} (adaptive Gaussian quadrature). A closed-form antiderivative
exists but requires case analysis; numerical evaluation is safer and
sufficiently fast for $N\leq 390$.

\subsection{Parallelisation}

Lambda-grid search and dataset sweep use \texttt{ProcessPoolExecutor} (process-level
parallelism) rather than \texttt{ThreadPoolExecutor}, because the Python GIL prevents
threads from achieving true CPU parallelism for numerical work~\needscite{Python GIL and multiprocessing, e.g.\ Python documentation}.
Worker functions are defined at module level to support cross-platform pickling
(required on Windows).

\subsection{Random Number Generation}

Path generation uses \texttt{numpy.random.default\_rng} (PCG64 generator), which is
thread-safe, statistically superior to the legacy Mersenne Twister, and reproducible
from integer seeds~\needscite{PCG random number generators, e.g.\ O'Neill (2014)}.
A master RNG generates per-path sub-seeds via \texttt{rng.integers}, ensuring reproducibility
without correlation between paths.

\subsection{Paired Comparison Design}

The comparator evaluates both trade schedules against the same Heston path, making the
IS difference purely attributable to the scheduling choice. An earlier version used
independently drawn paths for each strategy, which invalidated the paired $t$-test and
Wilcoxon test. The fix was identified when the paired test showed zero-variance differences
across all 15 datasets, confirming numerical collapse of the corrected schedule to AC.

\printbibliography

\end{document}
